\documentclass[a4paper, serif, 10pt]{moderncv}

%% moderncv
% color and style
\moderncvcolor{black}
\moderncvstyle{banking}

%% page layout/margins
\usepackage[top=2.5cm, bottom=2.5cm, left=1.5cm, right=1.5cm]{geometry}

\nopagenumbers{}

%% Babel config for Indonesian (Bahasa Indonesia) language
% ref:
% - https://latex3.github.io/babel/guides/locale-indonesian.html
\usepackage[indonesian]{babel}

%% fontspec
\usepackage{fontspec}

\IfFontExistsTF{Linux Libertine}{
  \setmainfont[Ligatures={TeX, Common}, Numbers=OldStyle]{Linux Libertine}
}{}
\IfFontExistsTF{Linux Libertine O}{
  \setmainfont[Ligatures={TeX, Common}, Numbers=OldStyle]{Linux Libertine O}
}{}

%% CTeX
% CJK support, used to show CJK characters in the resume
%
% - heading=false: disable ctex heading
% - punct=kaiming: use kaiming punctuations styles for CJK
% - scheme=plain: use plain scheme, do not override \normalsize font size
% - space=auto: space settings for CJK characters
%
% CJK fonts are resolved per locale language, e.g. Noto Serif/Sans CJK JP
% for Japanese resumes, with fallback to the next region if not installed
%
% ref:
% - http://ctan.mirrorcatalogs.com/language/chinese/ctex/ctex.pdf
\usepackage[UTF8, heading=false, punct=kaiming, scheme=plain, space=auto]{ctex}

\IfFontExistsTF{Noto Serif CJK SC}{
  \setCJKmainfont{Noto Serif CJK SC}
}{}
\IfFontExistsTF{Noto Sans CJK SC}{
  \setCJKsansfont{Noto Sans CJK SC}
}{}

%% line spacing
\usepackage{setspace}
\setstretch{1.125}

%% URL styling - use normal text instead of monospace
\urlstyle{same}

%% Patch \httplink and \httpslink to handle full URLs
% The original moderncv macros always prepend http:// or https:// to URLs.
% This causes issues when the URL already has a protocol (e.g., https://example.com)
% resulting in malformed URLs like https://https://example.com.
% This patch checks if the URL already contains "://" and uses it directly if so.
% Note: We use \str_set:Nx (x-expansion) to expand macros like \@homepage first.
\ExplSyntaxOn
\str_new:N \g_yamlresume_url_str

\RenewDocumentCommand{\httpslink}{O{}m}{
  \str_set:Nx \g_yamlresume_url_str {#2}
  \str_if_in:NnTF \g_yamlresume_url_str {://}
    { \url{#2} }
    { \url{https://#2} }
}

\RenewDocumentCommand{\httplink}{O{}m}{
  \str_set:Nx \g_yamlresume_url_str {#2}
  \str_if_in:NnTF \g_yamlresume_url_str {://}
    { \url{#2} }
    { \url{http://#2} }
}
\ExplSyntaxOff

\name{Arif Prasetyo}{}
\title{Insinyur Mesin Senior | Spesialis Manufaktur \& Otomasi}
\phone[mobile]{+62 812-3456-7890}
\email{arif.prasetyo.engineer@gmail.com}
\homepage{https://www.arifprasetyo.id}

\address{Jl. Cikini Raya No. 42, Menteng, Jakarta Pusat, DKI Jakarta, Indonesia, 10330}{}{}

\extrainfo{{\small \faLinkedin}\ \href{https://www.linkedin.com/in/arif-prasetyo-me}{@arif-prasetyo-me} {} {} {} • {} {} {} 
{\small \faGithub}\ \href{https://github.com/arifprasetyo-me}{@arifprasetyo-me}}

\begin{document}

\maketitle

\section{Informasi Dasar}

\cvline{}{Insinyur mesin dengan pengalaman lebih dari 8 tahun di industri manufaktur otomotif dan kedirgantaraan, berfokus pada perancangan mesin, optimasi proses produksi, dan implementasi otomasi.
%
\begin{itemize}
\item Menguasai CAD/CAE, analisis struktural, dan metodologi Lean Six Sigma untuk meningkatkan efisiensi lini produksi.
\item Berpengalaman memimpin proyek perbaikan proses lintas fungsi dengan penghematan biaya terukur dan peningkatan kualitas produk.
\item Terbiasa bekerja dengan standar ISO 9001, APQP, PPAP, serta keselamatan kerja K3 di lingkungan manufaktur.
\end{itemize}}
\section{Pendidikan}

\cventry{Agu 2013–Agu 2017}
        {Sarjana, Teknik Mesin, Nilai: 3.72}
        {Institut Teknologi Bandung}
        {\url{https://www.itb.ac.id}}
        {}
        {\begin{itemize}
\item Menyelesaikan tugas akhir tentang optimasi desain rangka mesin press menggunakan analisis elemen hingga.
\item Aktif dalam tim robotika kampus dan kompetisi perancangan mesin tingkat nasional.
\item Memperoleh pemahaman kuat dalam perancangan komponen, pemilihan material, dan proses produksi.
\end{itemize}
\textbf{Mata Kuliah}: Mekanika Fluida, Termodinamika, Perpindahan Panas, Kinematika dan Dinamika, Perancangan Mesin, Proses Manufaktur, Material Teknik, Getaran Mekanis}
\section{Pengalaman Kerja}

\cventry{Mar 2021–Sekarang}
        {Insinyur Mesin Senior}
        {PT Astra Otoparts Tbk}
        {\url{https://www.astra-otoparts.com}}
        {}
        {\begin{itemize}
\item Memimpin perancangan dan perbaikan fixture serta tooling untuk lini produksi komponen rem dengan target efisiensi 12\%.
\item Mengimplementasikan proyek otomasi berbasis PLC dan robotik untuk mengurangi waktu siklus dan kesalahan operator.
\item Melakukan analisis kegagalan komponen menggunakan metode root cause analysis dan finite element analysis.
\item Mengoordinasikan tim lintas fungsi dalam validasi proses, APQP, PPAP, dan audit standar ISO 9001.
\item Membimbing tiga insinyur junior dalam praktik desain untuk manufaktur dan keselamatan kerja.
\end{itemize}
\textbf{Kata Kunci}: Desain Manufaktur, Otomasi Produksi, Lean Manufacturing, CAD/CAE, ISO 9001}

\cventry{Jun 2018–Feb 2021}
        {Insinyur Proses Produksi}
        {PT Komatsu Indonesia}
        {\url{https://www.komatsu.com}}
        {}
        {\begin{itemize}
\item Menganalisis dan mengoptimalkan proses permesinan komponen alat berat untuk meningkatkan throughput sebesar 15\%.
\item Mengembangkan program CNC dan standar operasional untuk mengurangi scrap rate pada komponen kritis.
\item Memimpin kaizen mingguan bersama tim produksi dan maintenance untuk mengurangi downtime mesin.
\item Menyusun laporan analisis kapasitas, time study, dan rekomendasi investasi peralatan produksi.
\end{itemize}
\textbf{Kata Kunci}: CNC Programming, Lean Manufacturing, Six Sigma, Time Study, Perbaikan Berkelanjutan}

\cventry{Sep 2017–Mei 2018}
        {Insinyur Muda Perawatan}
        {PT Dirgantara Indonesia}
        {\url{https://www.indonesian-aerospace.com}}
        {}
        {\begin{itemize}
\item Mendukung perawatan preventif dan korektif komponen struktur pesawat sesuai manual dan standar keselamatan.
\item Melakukan inspeksi dimensional dan dokumentasi hasil perawatan untuk audit kualitas internal.
\item Berkolaborasi dengan teknisi senior dalam analisis kerusakan dan penyusunan prosedur perbaikan.
\end{itemize}
\textbf{Kata Kunci}: Maintenance, Inspeksi Kualitas, Struktur Pesawat, K3}
\section{Bahasa}

\cvline{Indonesia}{Bahasa Ibu / Bilingual \hfill \textbf{Kata Kunci}: Bahasa ibu}
\cvline{Inggris}{Tingkat Profesional Penuh \hfill \textbf{Kata Kunci}: TOEFL iBT 105, IELTS 7.0}
\section{Keahlian}

\cvline{Desain Mekanik}{Ahli \hfill \textbf{Kata Kunci}: SolidWorks, CATIA, AutoCAD, GD\&T, Finite Element Analysis}
\cvline{Manufaktur \& Produksi}{Mahir \hfill \textbf{Kata Kunci}: Lean Manufacturing, Six Sigma, CNC Programming, APQP, PPAP}
\cvline{Analisis Teknik}{Mahir \hfill \textbf{Kata Kunci}: ANSYS, MATLAB, Thermofluids, Root Cause Analysis, DOE}
\cvline{Manajemen Proyek}{Menengah \hfill \textbf{Kata Kunci}: MS Project, Scrum, Budgeting, Vendor Management}
\section{Penghargaan}

\cventry{Nov 2023}
        {PT Astra Otoparts Tbk}
        {Penghargaan Inovasi Manufaktur Terbaik}
        {}
        {}
        {Diberikan atas keberhasilan memimpin proyek otomasi lini perakitan komponen rem yang meningkatkan produktivitas dan menurunkan biaya operasional.}

\cventry{Des 2020}
        {PT Komatsu Indonesia}
        {Karyawan Terbaik Tahun 2020}
        {}
        {}
        {Diakui karena kontribusi dalam optimasi proses permesinan dan penurunan scrap rate pada komponen kritis.}
\section{Sertifikat}

\cventry{Mei 2019}
        {Dassault Systèmes}
        {Certified SolidWorks Professional}
        {\url{https://www.solidworks.com/certification}}
        {}
        {}

\cventry{Okt 2020}
        {ASQ}
        {Six Sigma Green Belt}
        {\url{https://asq.org/cert/six-sigma-green-belt}}
        {}
        {}

\cventry{Feb 2018}
        {Kementerian Ketenagakerjaan Republik Indonesia}
        {Ahli K3 Umum}
        {\url{https://kemnaker.go.id}}
        {}
        {}
\section{Publikasi}

\cventry{Apr 2022}
        {Analisis Pengurangan Waktu Setup Mesin CNC pada Produksi Komponen Otomotif}
        {Jurnal Teknik Mesin Indonesia}
        {\url{https://doi.org/10.1234/jtmi.2022.012}}
        {}
        {\begin{itemize}
\item Membahas pendekatan lean manufacturing dan standardisasi tooling untuk mengurangi waktu setup mesin CNC.
\item Menyajikan data time study dan dampak pengurangan waktu setup terhadap kapasitas produksi.
\item Merekomendasikan praktik terbaik untuk implementasi perbaikan berkelanjutan di lantai produksi.
\end{itemize}}
\section{Proyek}

\cventry{Jan 2022–Jun 2023}
        {Proyek otomasi lini perakitan komponen rem yang mencakup desain fixture, integrasi PLC, dan validasi proses produksi.
}
        {Sistem Otomasi Lini Perakitan Komponen Rem}
        {}
        {}
        {\begin{itemize}
\item Merancang fixture dan mekanisme feeding untuk mengurangi waktu siklus perakitan.
\item Mengintegrasikan sistem PLC dan sensor untuk meningkatkan konsistensi kualitas produk.
\item Menghasilkan penghematan biaya tahunan dan penurunan cacat produk.
\end{itemize}
\textbf{Kata Kunci}: Otomasi Lini Produksi, PLC, Fixture Design, Validasi Proses}

\cventry{Feb 2020–Des 2020}
        {Proyek rekayasa ulang rangka mesin press untuk meningkatkan kekakuan struktur dan mengurangi berat material.
}
        {Optimasi Desain Rangka Mesin Press 120 Ton}
        {}
        {}
        {\begin{itemize}
\item Melakukan analisis elemen hingga untuk mengevaluasi tegangan dan deformasi rangka.
\item Mengusulkan perubahan geometri dan material yang menurunkan berat tanpa mengorbankan kekuatan.
\item Memvalidasi desain melalui pengujian prototipe dan simulasi beban kerja.
\end{itemize}
\textbf{Kata Kunci}: Desain Rangka, Optimasi Struktur, Finite Element Analysis, Pengurangan Biaya}
\section{Minat}

\cvline{Olahraga}{Futsal, Bersepeda, Panjat Tebing}
\cvline{Musik}{Gitar, Keyboard}
\cvline{Otomotif}{Restorasi Motor Klasik, Touring}
\section{Kegiatan Sukarela}

\cventry{Jan 2021–Sekarang}
        {Mentor Teknis}
        {Persatuan Insinyur Indonesia Komisariat Jakarta}
        {\url{https://pii.or.id}}
        {}
        {\begin{itemize}
\item Membimbing mahasiswa teknik mesin dalam perancangan proyek akhir dan persiapan karier industri.
\item Menjadi narasumber lokakarya tentang CAD, manufaktur lean, dan keselamatan kerja.
\item Mendukung kegiatan komunitas insinyur muda untuk meningkatkan kompetensi teknis dan profesionalisme.
\end{itemize}}
    
\end{document}